\section{Resultados}
\subsection{Evaluación en DRIVE (in-distribution)}
La Tabla~\ref{tab:drive} resume el desempeño del modelo principal en DRIVE.

\begin{table}[h]
\centering
\caption{Resultados en DRIVE (completar con valores de ejecución).}
\label{tab:drive}
\begin{tabular}{lcccc}
\toprule
Modelo & Sensibilidad & Especificidad & F1 & AUC-ROC \\
\midrule
Attention U-Net + pérdida combinada & \textbf{--} & \textbf{--} & \textbf{--} & \textbf{--} \\
\bottomrule
\end{tabular}
\end{table}

\subsection{Generalización DRIVE $\rightarrow$ CHASE\_DB1}
La Tabla~\ref{tab:cross} reporta la evaluación fuera de distribución y la brecha respecto a DRIVE.

\begin{table}[h]
\centering
\caption{Generalización cruzada y brecha de rendimiento.}
\label{tab:cross}
\begin{tabular}{lcccc}
\toprule
Dataset & Sensibilidad & Especificidad & F1 & AUC-ROC \\
\midrule
DRIVE & -- & -- & -- & -- \\
CHASE\_DB1 & -- & -- & -- & -- \\
Brecha ($\Delta$ DRIVE-CHASE) & -- & -- & -- & -- \\
\bottomrule
\end{tabular}
\end{table}

\subsection{Efecto de la adaptación de dominio}
La Tabla~\ref{tab:adapt} muestra el efecto de la estrategia de adaptación por preprocesamiento en CHASE\_DB1.

\begin{table}[h]
\centering
\caption{Impacto de la adaptación (antes vs después).}
\label{tab:adapt}
\begin{tabular}{lccc}
\toprule
Estrategia & F1 & Sensibilidad & AUC-ROC \\
\midrule
Sin adaptación & -- & -- & -- \\
CLAHE + normalización & -- & -- & -- \\
CLAHE + normalización + TTA & -- & -- & -- \\
\bottomrule
\end{tabular}
\end{table}

\subsection{Análisis cualitativo de fallos}
De forma consistente, los mayores errores se concentran en capilares finos y regiones de bajo contraste. En cambio, arterias y venas de mayor calibre muestran contornos más estables. Este patrón sugiere una limitación estructural: la pérdida de detalle inducida por downsampling afecta más a estructuras delgadas, incluso cuando existen skip connections.

Se recomienda incluir en la versión final una figura comparativa con: (a) imagen original, (b) máscara real, (c) predicción, (d) mapa de error; mostrando al menos casos de éxito y fallo en capilares.
